% !TeX program = xelatex
% 使用 XeLaTeX 编译两次，以更新目录和页码。
\documentclass[10pt,a4paper,fontset=none]{ctexart}
\usepackage[margin=19mm,top=17mm,bottom=18mm,footskip=8mm]{geometry}
\usepackage{amsmath,amssymb}
\usepackage{xcolor}
\usepackage{enumitem}
\usepackage{fvextra}
\usepackage{hyperref}

\definecolor{Ink}{HTML}{20252B}
\definecolor{Heading}{HTML}{233B52}
\definecolor{Rule}{HTML}{9CA9B4}
\hypersetup{
  hidelinks,
  pdftitle={算法竞赛模板},
  pdfdisplaydoctitle=true,
  bookmarksnumbered=true
}

\setmonofont{DejaVuSansMono.ttf}
\setCJKmainfont{gbsn00lp.ttf}[AutoFakeBold=2,ItalicFont=gkai00mp.ttf]
\setCJKmonofont{gkai00mp.ttf}
\setcounter{tocdepth}{2}
\setlength{\parindent}{0pt}
\setlength{\parskip}{2pt plus 1pt}
\setlist{nosep,leftmargin=1.7em,topsep=2pt}
\setlength{\emergencystretch}{1em}
\ctexset{
  section={
    format=\large\bfseries\color{Heading},
    beforeskip=1.7ex plus .3ex,
    afterskip=.8ex,
    aftername=\enspace
  },
  subsection={
    format=\normalsize\bfseries\color{Heading},
    beforeskip=1.3ex plus .3ex,
    afterskip=.5ex,
    aftername=\enspace
  },
  subsubsection={
    format=\small\bfseries,
    beforeskip=.9ex plus .2ex,
    afterskip=.3ex
  }
}
\newcommand{\CodeFont}{\fontsize{9.4pt}{10.2pt}\selectfont}
\newcommand{\CppInput}[1]{%
  \VerbatimInput[
    fontsize=\CodeFont,
    breaklines=true,
    breakanywhere=true,
    breakautoindent=false,
    breaksymbolleft={\(\hookrightarrow\)},
    breaksymbolright={},
    breakindent=.8em,
    tabsize=4,
    frame=leftline,
    framesep=3pt,
    rulecolor=\color{Rule},
    framerule=.35pt,
    xleftmargin=4pt,
    listparameters={\setlength{\topsep}{3pt}}
  ]{#1}%
}

\pagestyle{plain}
\begin{document}
\color{Ink}
{\color{Heading}\LARGE\bfseries 算法竞赛模板\par}
\vspace{3pt}
{\color{Rule}\hrule height .5pt}
{\small\setlength{\parskip}{0pt}\tableofcontents}
\vspace{1pt}
{\color{Rule}\hrule height .5pt}

\section{数论}

\subsection{扩展欧几里得}

\CppInput{code/number_theory/extended_gcd.cpp}

\subsection{中国剩余定理}

中国剩余定理（CRT）用于求解同余方程组。

模数互质（标准CRT）

设模数 $m_1, m_2, ..., m_k$ 两两互质。

\begin{itemize}
\item 令 $M = \prod m_i$
\item $M_i = M/m_i$
\end{itemize}

由于 $M_i$ 与 $m_i$ 互质，存在逆元 $t_i$ 使 $M_i t_i \equiv 1 \pmod{m_i}$。

\textbf{解为：}
\[
x \equiv \sum_{i=1}^{k} a_i M_i t_i \pmod{M}
\]

\textbf{验证：}
对模 $m_j$，除第 $j$ 项为 $a_j$ 外，其余项均含 $m_j$ 因子而消去。

\CppInput{code/number_theory/chinese_remainder_theorem.cpp}

\subsection{线性筛}

\CppInput{code/number_theory/linear_sieve.cpp}

\subsection{欧拉函数}

\CppInput{code/number_theory/euler_totient.cpp}

\subsection{一阶线性同余方程}

\subsubsection*{标准形式}
\[
ax \equiv b \pmod{m}
\]

求解 $x$，或判断无解。

\subsubsection*{求解步骤}

\subsubsection*{1. 化为不定方程}
\[
ax \equiv b \pmod{m} \iff ax + my = b
\]

\subsubsection*{2. 有解条件}

设 $d = \gcd(a, m)$，方程有解当且仅当：
\[
d \mid b
\]

\subsubsection*{3. 求特解}

用扩展欧几里得解 $ax_0 + my_0 = d$，两边乘 $\frac{b}{d}$：
\[
x_1 = x_0 \cdot \frac{b}{d}
\]

\subsubsection*{4. 求通解}

模 $m$ 意义下有 $d$ 个解：
\[
\begin{gathered} x \equiv x_1 + k \cdot \frac{m}{d} \pmod{m}, \\ k = 0, 1, \ldots, d-1 \end{gathered}
\]

\textbf{最小非负解：} $x_1 \bmod \frac{m}{d}$

\subsubsection*{模板代码}

\CppInput{code/number_theory/linear_congruence.cpp}

\subsection{分解质因数}

\CppInput{code/number_theory/prime_factorization.cpp}

\subsection{二次剩余}

\subsubsection*{定义}

对于奇素数 $p$ 和整数 $a$（$p \nmid a$），若存在 $x$ 使得：
\[
x^2 \equiv a \pmod{p}
\]

则称 $a$ 是模 $p$ 的\textbf{二次剩余}，否则为\textbf{非二次剩余}。

\subsubsection*{勒让德符号}
\[
\left(\frac{a}{p}\right) =
\begin{cases}
1  & a \text{ 是二次剩余} \\
-1 & a \text{ 是非二次剩余} \\
0  & p \mid a
\end{cases}
\]

\subsubsection*{欧拉判别法}
\[
\left(\frac{a}{p}\right) \equiv a^{\frac{p-1}{2}} \pmod{p}
\]

\subsubsection*{Cipolla 算法（求解二次剩余）}

\textbf{复杂度}：$O(\log p)$

\textbf{步骤}：
\begin{enumerate}
\item 若 $p=2$，$x \equiv a \pmod{2}$
\item 随机找 $r$ 使 $r^2 - a$ 是非二次剩余
\item 在扩域 $\mathbb{F}_p[\omega] = \{u + v\omega \mid \omega^2 = r^2 - a\}$ 中计算：
\[
x \equiv (r + \omega)^{\frac{p+1}{2}} \pmod{p}
\]

\end{enumerate}

\subsubsection*{模板代码}

\CppInput{code/number_theory/cipolla.cpp}

\subsection{质因数个数}

\CppInput{code/number_theory/prime_factor_counts.cpp}

\subsection{线性基}

\CppInput{code/number_theory/xor_linear_basis.cpp}

\subsection{组合数}

在素数模 $p$ 下，预处理 $0$ 到 $N$ 的阶乘与逆阶乘，要求 $N < p$。
例如 \texttt{Combination comb(100);} 后，\texttt{comb.C(5, 2)} 为 $10$。
预处理 $O(N+\log p)$，每次查询 $O(1)$。

\CppInput{code/number_theory/combination.cpp}

\section{图论}

\subsection{Tarjan}

\subsubsection*{算法思想}

对有向图进行强连通分量（SCC）缩点，缩点后得到 DAG。

\textbf{核心数组：}
\begin{itemize}
\item \texttt{dfn[u]}：DFS 序
\item \texttt{low[u]}：能回溯到的最早 dfn
\item \texttt{in\_stk[u]}：是否在栈中
\item \texttt{scc[u]}：所属 SCC 编号
\end{itemize}

\textbf{判定 SCC：} \texttt{dfn[u] == low[u]} 时，栈顶到 \texttt{u} 的所有点构成一个 SCC。

\subsubsection*{模板代码}

\CppInput{code/graph/tarjan_scc.cpp}

\subsection{网络流Dinic解法}

\CppInput{code/graph/dinic.cpp}

\subsection{树的直径树形DP}

\CppInput{code/graph/tree_diameter.cpp}

\subsection{费用流Dinic+SPFA}

\CppInput{code/graph/min_cost_max_flow.cpp}

\subsection{DFS序求LCA}

\CppInput{code/graph/dfs_order_lca.cpp}

\section{字符串}

\subsection{KMP}

\CppInput{code/string/kmp.cpp}

\subsection{Trie}

支持小写英文字母，可统计完整字符串和前缀的出现次数。
例如插入 \texttt{ab} 和 \texttt{abc} 后，\texttt{count("ab")} 为 $1$，
\texttt{count\_prefix("ab")} 为 $2$。单次操作为 $O(|s|)$。

\CppInput{code/string/trie.cpp}

\section{数据结构}

\subsection{线段树}

\CppInput{code/data_structure/segment_tree.cpp}

\subsection{左偏树可并堆（小根堆）}

\CppInput{code/data_structure/leftist_heap.cpp}

\subsection{单调栈}

求每个位置两侧最近的严格更小元素下标，下标从 $0$ 开始，不存在时为 $-1$。
两处弹栈条件的 \texttt{>=} 改为 \texttt{<=}，即可求严格更大的元素。复杂度 $O(n)$。

\CppInput{code/data_structure/monotonic_stack.cpp}

\subsection{单调队列}

求每个长度为 $k$ 的窗口的最值，默认求最大值。传入 \texttt{false} 求最小值。
例如对 $\{3,1,2\}$ 取 $k=2$，得到最大值 $\{3,2\}$、最小值 $\{1,1\}$。复杂度 $O(n)$。

\CppInput{code/data_structure/monotonic_queue.cpp}

\section{动态规划}

\subsection{最长上升子序列}

用二分维护各长度子序列的最小末尾，通过前驱下标恢复一组严格上升子序列。
例如 $\{3,1,1,2\}$ 可得到 $\{1,2\}$。时间 $O(n\log n)$，额外空间 $O(n)$。

\CppInput{code/dynamic_programming/lis.cpp}

\section{线性代数}

\subsection{高斯消元}

求解 $n$ 个方程、$m$ 个未知数的实数线性方程组。传入 $n \times (m+1)$ 的增广矩阵，
在每列选取绝对值最大的主元，以 $10^{-9}$ 为判零阈值。返回值 $0$、$1$、$2$ 分别表示唯一解、无穷多解、无解。
有多解时，返回自由变量取零的一组解。时间 $O(nm\min(n,m))$。

\CppInput{code/linear_algebra/gaussian_elimination.cpp}

\end{document}
